% =============================================================
%  ch16_projekt2.tex  --  Projekt 2: Tischtennisroboter
%  Gliederung: Vorueberlegungen (Architektur, Mechanik, CAD/URDF) und
%  Vorgehensweise (zuerst Wahrnehmung simulieren, dann Roboter).
% =============================================================

\chapter{Projekt 2 -- Tischtennisroboter}
\label{ch:projekt2}

\textbf{Ziel:} Ballflug per Stereokamera + Kalman schätzen, Roboter fährt an
die Abschlagposition und schätzt Abschlagzeitpunkt und -trajektorie.

\begin{infobox}[title={Grundlagen aus den Einstiegsprojekten werden vorausgesetzt}]
Die \emph{Bausteine} sind in den Einstiegsprojekten eingeführt und werden hier
nur referenziert: \textbf{Stereo-Triangulation und Kalman-Tracking} in
Projekt~0.B (Kap.~\ref{ch:projekt0b}); \textbf{\code{ros2\_control}, URDF/Xacro,
Topics, Services und Actions} in Projekt~0.A (Kap.~\ref{ch:projekt0a}). Dieses
Kapitel gliedert sich in \textbf{Vorüberlegungen} (Architektur, Mechanik,
Simulierbarkeit, CAD/URDF) und \textbf{Vorgehensweise} (was wir in welcher
Reihenfolge umsetzen).
\end{infobox}

% =============================================================
\section{Vorüberlegungen}
Bevor gebaut wird, klären wir die Software-Bausteine, wo gerechnet wird, den
mechanischen Aufbau und seine Grenzen in der Simulation sowie die nebenläufige
Ablauflogik.

\subsection{ROS\,2-Bausteine}
\begin{longtable}{@{}p{2.6cm}p{12cm}@{}}
\toprule
\textbf{Typ} & \textbf{Konkret} \\
\midrule
Nodes   & \code{stereo\_camera} (Treiber) $\cdot$
          \code{stereo\_image\_proc} (Disparität/Tiefe) $\cdot$
          \code{ball\_detector} (Farb-/Blob-Erkennung $\rightarrow$ 3D-Position) $\cdot$
          \code{ball\_tracker} (EKF, Flugvorhersage) $\cdot$
          \code{strike\_planner} (Abfangpunkt + Zeit + Schlagbahn) $\cdot$
          \code{robot\_controller} (Projekt-1-Stack) \\
Topics  & \code{/left/image\_raw}, \code{/right/image\_raw} $\cdot$
          \code{/points2} $\cdot$
          \code{/ball/position} $\cdot$
          \code{/ball/predicted\_trajectory} $\cdot$
          \code{/robot/target\_pose} \\
Actions & \code{FollowJointTrajectory} (Schlagbewegung) \\
Parameter & Kamerakalibrierung, Ballfarbe/-radius, EKF-$Q$/$R$, Tischgeometrie \\
\bottomrule
\end{longtable}

\subsection{Architektur und Wahrnehmungskette}
\begin{center}
\begin{tikzpicture}[node distance=5mm and 11mm]
\node[node] (cam) {/stereo\_camera};
\node[node,right=of cam] (sip) {/stereo\_image\_proc};
\node[node,right=of sip] (det) {/ball\_detector};
\node[node,right=of det] (trk) {/ball\_tracker\\\scriptsize (EKF)};
\node[node,below=9mm of trk] (plan) {/strike\_planner};
\node[node,left=of plan] (rc) {/robot\_controller};
\draw[flow] (cam) -- node[topic,fill=white,inner sep=1pt]{\scriptsize L/R img} (sip);
\draw[flow] (sip) -- node[topic,fill=white,inner sep=1pt]{\scriptsize points2} (det);
\draw[flow] (det) -- node[topic,fill=white,inner sep=1pt]{\scriptsize ball/pos} (trk);
\draw[flow] (trk) -- node[topic,fill=white,inner sep=1pt]{\scriptsize pred traj} (plan);
\draw[flow] (plan) -- node[topic,fill=white,inner sep=1pt]{\scriptsize target pose} (rc);
\end{tikzpicture}
\captionof{figure}{Wahrnehmungskette von der Stereokamera über den
EKF bis zum Schlagplaner und Roboter.}
\end{center}
\noindent\textbf{Kernidee:} Der EKF (Kapitel~\ref{ch:kalman}) bekommt die
verrauschten 3D-Ballpositionen und extrapoliert per reiner Prädiktion bis zur
Schlagebene; daraus folgen Abfangpunkt und -zeit. \textbf{Hauptsächlich reales
Problem:} Latenz und Timing -- in der Simulation triviale Zeit, real kritisch.

\subsection{Kamera-Architektur -- wo wird gerechnet?}
Lohnt es sich, die Stereokamera an den ESP32 / Raspberry Pi Zero zu hängen und
die Rohdaten per ROS-Service oder einfachem UDP/WLAN weiterzuleiten -- oder
entstehen Zeit-/Rechenleistungsprobleme?

\subsubsection{Die Bandbreitenrechnung entscheidet}
Die synchronisierte Stereokamera liefert ein side-by-side-Bild
(z.\,B.\ $2\times1280\times720$). Was das pro Sekunde bedeutet:
\begin{longtable}{@{}p{6.2cm}rr@{}}
\toprule
\textbf{Datenstrom} & \textbf{pro Frame} & \textbf{bei 60\,fps} \\
\midrule
Stereo roh, YUYV (2\,B/px), $2560\times720$ & 3{,}69\,MB & $\approx$\,1769\,Mbit/s \\
Stereo MJPEG-komprimiert ($\approx$\,1:10)   & 0{,}37\,MB & $\approx$\,177\,Mbit/s \\
\textbf{Nur 3D-Ballposition} (Header + 3 Floats) & $\approx$\,80\,B & $\approx$\,0{,}1\,Mbit/s \\
\bottomrule
\end{longtable}
\noindent Dem gegenüber die realistischen Linkkapazitäten:
\begin{longtable}{@{}p{6.2cm}r@{}}
\toprule
\textbf{Verbindung} & \textbf{real nutzbar} \\
\midrule
ESP32 WLAN (TCP, real)         & $\approx$\,10--30\,Mbit/s \\
Pi Zero (2{,}4\,GHz WLAN)     & einige zehn Mbit/s, hoher Jitter \\
Pi\,4 WLAN (Dualband)          & $\approx$\,100--300\,Mbit/s, mit Jitter \\
Pi\,4 Gigabit-Ethernet (Kabel) & 1000\,Mbit/s, deterministisch \\
USB\,2 / USB\,3               & 480 / 5000\,Mbit/s \\
\bottomrule
\end{longtable}

\subsubsection{Verdikt pro Plattform}
\begin{description}[style=nextline,leftmargin=1.4em]
\item[ESP32 -- nein.] Zwei harte Gründe: (1)~Der ESP32 ist \emph{kein}
  USB-Host für UVC-Kameras. (2)~177--1769\,Mbit/s liegen um Größenordnungen
  über seinem WLAN-Durchsatz. Der ESP32 gehört auf die
  \textbf{Echtzeit-Motorebene}, nicht an die Kamera.
\item[Pi Zero -- technisch möglich, aber schlechte Idee.] Der Zero\,2\,W kann
  mit OTG-Adapter eine UVC-Kamera einlesen, aber: schwache CPU, nur
  2{,}4-GHz-WLAN, und er wäre allein mit dem Datenschaufeln gesättigt. Für
  einen \emph{latenzkritischen} Balltracker ist WLAN-Jitter Gift.
\item[Pi\,4 -- ja, und das ist der richtige Ansatz.] Der Pi\,4 hat USB\,3,
  Gigabit-Ethernet und eine brauchbare Quad-Core-CPU. Er kann die Kamera
  einlesen \emph{und die Bildverarbeitung lokal machen}.
\end{description}
\begin{infobox}[title={Das richtige Muster: am Rand rechnen, nur das Ergebnis senden}]
Schiebe nicht die Pixel durchs Netz, sondern die \textbf{Erkenntnis}. Lass den
Pi\,4 (direkt an der Kamera) die Stereo-Erkennung machen und nur die
\textbf{3D-Ballposition} publizieren -- das sind $\approx$\,0{,}1\,Mbit/s,
trivial über jedes Netz. ROS\,2/DDS transportiert das Topic
netzwerktransparent zum PC. Für deterministische Latenz: \textbf{Pi\,4 per
Kabel-Gigabit-Ethernet} an den PC.
\end{infobox}

\subsubsection{Sparse statt dense: der Ball braucht keine volle Tiefenkarte}
Eine dichte Disparitätskarte (SGBM über das ganze Bild) ist teuer. Für die
Ballverfolgung genügt \textbf{sparse Triangulation}: den Ball (heller Blob) in
\emph{beiden} rektifizierten Bildern detektieren und über $Z=fb/d$ nur die Tiefe
\emph{dieses einen Punktes} bestimmen. Genau diese Pixel-zu-3D-Triangulation
haben wir in Projekt~0.B (Kap.~\ref{ch:projekt0b}) implementiert; sie läuft auf
dem Pi\,4 locker bei 60--120\,fps.

\subsubsection{ROS-Service, UDP oder Topic?}
\begin{itemize}
\item \textbf{ROS-Service: falsch} für Streaming. Services sind synchrones
  Request/Response (blockierend, one-shot) -- für kontinuierliche Ströme
  ungeeignet.
\item \textbf{ROS-Topic: richtig.} Ergebnis (Ballposition) als
  \code{RELIABLE}-Topic; falls doch Bilder, \code{image\_transport} mit
  \code{compressed} und QoS \code{BEST\_EFFORT}.
\item \textbf{Rohes UDP/GStreamer: nur als letzter Ausweg} -- du verlierst die
  ROS-Integration. Für berechnete Ergebnisse unnötig.
\end{itemize}
\begin{center}
\begin{tikzpicture}[node distance=6mm and 14mm]
\node[node,fill=accent2!8,draw=accent2] (cam) {Stereokamera\\\scriptsize USB};
\node[node,right=of cam] (pi)
  {Raspberry Pi\,4\\\scriptsize Stereo + Balldetektion\\\scriptsize (sparse Triangulation)};
\node[node,right=22mm of pi] (pc) {Externer PC\\\scriptsize EKF + Schlagplaner};
\node[node,below=8mm of pc] (esp) {ESP32 (micro-ROS)\\\scriptsize Achsen / Schrittmotoren};
\draw[flow] (cam) -- node[topic,fill=white,inner sep=1pt]{\scriptsize USB3} (pi);
\draw[flow] (pi)  -- node[topic,fill=white,inner sep=1pt]{\scriptsize /ball/position (GbE)} (pc);
\draw[flow] (pc)  -- node[topic,fill=white,inner sep=1pt]{\scriptsize Soll-Trajektorie} (esp);
\end{tikzpicture}
\captionof{figure}{Empfohlene Architektur: schwere Pixelarbeit am Pi\,4 nahe der
Kamera, nur die destillierte Ballposition wandert (kabelgebunden) zum PC.}
\end{center}

\subsection{Mechanischer Aufbau: Portalsystem (Gantry)}
\subsubsection{Grundkonzept}
Statt eines Delta-Roboters ein \textbf{kartesisches Portalsystem}: große
Linearachsen positionieren den Schläger grob im Raum, ein kleiner Arm am Ende
führt die Schlagbewegung aus und stellt den Schlägerwinkel. Bezug
Tischtennisplatte (ITTF): $2{,}74\,\text{m} \times 1{,}525\,\text{m}$, Höhe
$0{,}76\,\text{m}$, Netz $15{,}25\,\text{cm}$. Die \textbf{X-Achse} überspannt
die Plattenbreite ($1{,}525\,\text{m}$).

\subsubsection{Geschachtelte Achsen}
\begin{itemize}
\item \textbf{X-Achse} (Plattenbreite, horizontal): zwei parallele
  Führungsstäbe, Zahnriemenantrieb, X-Schlitten. \emph{Glatte Hohlstahl-Stäbe +
  GT2-Zahnriemen}, leichter Schlitten.
\item \textbf{Z-Achse} (Höhe, vertikal): nach demselben Prinzip, montiert
  \emph{im} X-Schlitten. Der Abstand der X-Führungsstäbe definiert den
  \textbf{Hub der Z-Achse}.
\item \textbf{Y-Achse} (Tiefe): entweder starr (fester Ausleger) oder als
  dritte Linearachse auf dem Z-Schlitten.
\item \textbf{Schlagarm}: auf dem Y/Z-Schlitten ein kleiner Roboterarm mit
  Schlägereinspannung.
\end{itemize}
\begin{warnbox}[title={Zahnriemen, nicht Keilriemen}]
Für \emph{Positionierung} brauchst du einen \textbf{Zahnriemen} (GT2),
keinen Keilriemen. Ein Keilriemen überträgt per Reibschluss und \emph{schlupft}
-- damit verlierst du Positionsgenauigkeit. Zahnriemen sind formschlüssig und
der Standard in 3D-Druckern/Plottern.
\end{warnbox}

\subsubsection{Antriebstopologie und Hardware-Anschluss}
\begin{longtable}{@{}p{3.2cm}p{11.4cm}@{}}
\toprule
\textbf{Element} & \textbf{Umsetzung} \\
\midrule
X-/Z-/Y-Motoren  & NEMA-17, möglichst \emph{rahmenfest} (Riemen treibt Schlitten) \\
Treiber          & TMC2209 (UART) je Achse, StallGuard als Last-/Endlagenhinweis \\
Endschalter      & optische/mechanische Endstops je Achse zum Homing \\
Strom/Spannung   & INA226-Modul je Treiberzweig (I\textsuperscript{2}C), optional \\
Controller       & ESP32 (micro-ROS): bündelt STEP/DIR bzw.\ UART, liest Endstops + INA226 \\
Versorgung       & 24\,V Netzteil für Motoren, getrennte 5\,V-Logik; Sternmasse \\
Schlagarm-Motoren & klein, idealerweise basisnah (Riemen-/Seilzug zu distalen Gelenken) \\
\bottomrule
\end{longtable}
\begin{center}
\begin{tikzpicture}[scale=1.0,every node/.style={font=\scriptsize}]
\draw[thick,accent] (0,0) rectangle (7,4.2);
\node[accent] at (3.5,4.5) {Rahmen (rahmenfest, X-Motor + Idler)};
\draw[gray,line width=2pt] (0.5,3.6)--(6.5,3.6);
\draw[gray,line width=2pt] (0.5,2.4)--(6.5,2.4);
\node[gray] at (-0.2,3.0) {X};
\draw[fill=accent!12,draw=accent] (2.6,2.2)rectangle(4.4,3.8);
\node[accent] at (3.5,3.95) {X-Schlitten};
\draw[accent2,line width=1.5pt] (2.9,2.3)--(2.9,3.7);
\draw[accent2,line width=1.5pt] (4.1,2.3)--(4.1,3.7);
\node[accent2] at (3.5,3.0) {Z};
\draw[fill=accent2!12,draw=accent2] (2.95,2.55)rectangle(4.05,3.0);
\node[accent2] at (3.5,2.78) {Z-Schlitten};
\draw[line width=1.4pt,black] (3.5,2.55)--(3.5,1.0);
\draw[fill=black!8] (3.2,0.6) rectangle (3.8,1.0);
\node at (3.5,0.8) {Arm};
\draw[fill=accent!20,draw=accent] (3.35,0.2) rectangle (3.65,0.6);
\node[rotate=90] at (3.5,0.4) {\tiny Schl\"ager};
\draw[<->,accent2] (4.0,2.55)--(4.0,0.4);
\node[accent2,right] at (4.05,1.5) {Hebelarm $L$};
\end{tikzpicture}
\captionof{figure}{Frontblick: X-Schlitten trägt die im Rod-Abstand
geschachtelte Z-Achse; der Schlagarm hängt als Ausleger (Hebelarm $L$).}
\end{center}

\subsection{Das Hebelarm- und Schwingungsproblem}
Der Arm hängt mit Abstand $L$ unter der X-Achse. Bei harten Beschleunigungen
zur Zielposition entstehen Reaktionskräfte und -momente, die den Ausleger zum
Schwingen anregen.

\subsubsection{Die Physik}
Beschleunigt der Schlitten mit $a$, wirkt auf die bewegte Masse $m$ die
Trägheitskraft $F=m\,a$. Über den Hebelarm $L$ ergibt das ein \textbf{Biegemoment}
\begin{equation}
M = F\cdot L = m\,a\,L.
\end{equation}
Der Ausleger verhält sich wie ein einseitig eingespannter Balken (Kragträger)
mit Steifigkeit $k=\dfrac{3EI}{L^3}$ und erster \textbf{Eigenfrequenz}
\begin{equation}
f_1 = \frac{1}{2\pi}\sqrt{\frac{k}{m_{\text{eff}}}}
    = \frac{1}{2\pi}\sqrt{\frac{3EI}{m_{\text{eff}}\,L^3}}.
\end{equation}
$f_1$ fällt mit $L^{3/2}$ -- ein längerer Ausleger ist drastisch weicher.
Anregung nahe $f_1$ führt zu Resonanz und Nachschwingen, was die
Schlaggenauigkeit ruiniert.

\subsubsection{Gegenmaßnahmen}
\begin{enumerate}
\item \textbf{Bewegte Masse minimieren:} Motoren rahmenfest (Riemen-/Seilzug,
  CoreXY-Logik), Arm aus Aluminium/CFK, Schwerpunkt nah an der Achse.
\item \textbf{Steifigkeit maximieren:} große Schienenquerschnitte (gestützte
  Linearschienen), Doppelführung, kurzer Ausleger, triangulierter Rahmen.
\item \textbf{Ruckbegrenzte Profile (S-Kurve):} sanftes Beschleunigen regt hohe
  Frequenzen weniger an.
\item \textbf{Input Shaping / Resonanzkompensation:} der Steuerbefehl wird so
  vorgeformt, dass die Anregung bei $f_1$ ausgelöscht wird -- wie Bambu/Klipper
  per ADXL345. Die wirksamste softwareseitige Maßnahme.
\item \textbf{Dämpfung:} Strukturdämpfung, abgestimmter Tilger, höhere
  Riemenvorspannung.
\end{enumerate}
\begin{infobox}[title=Anknüpfung an den FEM-Hintergrund]
Bestimme $f_1$ und die Modenformen per \textbf{Modalanalyse} (SolidWorks
Simulation / FEM). Das liefert die Frequenz, auf die du das Input Shaping
abstimmst, und zeigt, welche Strukturteile zu versteifen sind.
\end{infobox}

\subsection{Simulierbarkeit in Gazebo und RViz}
Die Sim-Real-Trennung über \code{ros2\_control} (Kap.~\ref{ch:simreal}) und der
Gazebo-Workflow sind aus 0.A/0.B bekannt -- hier nur die \emph{Grenzen}.
\subsubsection{Was Gazebo kann -- und was nicht}
\begin{longtable}{@{}p{7.4cm}p{7.2cm}@{}}
\toprule
\textbf{Gazebo bildet ab (Starrkörperdynamik)} & \textbf{Gazebo bildet \emph{nicht} ab} \\
\midrule
Anfahren der Zielposition (prism.\ + rot.\ Gelenke) & Strukturelle Biegung der Stäbe \\
Schlagbewegung des Arms      & Riemendehnung als kontinuierliche Elastizität \\
Gelenkmomente/-kräfte (auslesbar) & Modale Resonanz/Eigenschwingung \\
Dynamische Kopplung (Armträgheit $\leftrightarrow$ Portal) & Rahmenflattern, Mikrovibrationen \\
Ball-Schläger-Kontakt + Reaktionskräfte & exakte Spin-/Absprungphysik (nur genähert) \\
\bottomrule
\end{longtable}
\begin{warnbox}[title=Der Kernpunkt]
Gazebo behandelt Links als \textbf{starre Körper}. Die befürchtete Schwingung
kommt aber aus \emph{struktureller Nachgiebigkeit} (Stabbiegung, Riemendehnung)
-- starre Links biegen sich nicht, also zeigt Standard-Gazebo die Resonanz
\emph{nicht}.
\end{warnbox}
\subsubsection{Wege, die Nachgiebigkeit doch abzubilden}
\begin{itemize}
\item \textbf{Klumpen-Nachgiebigkeit in Gazebo:} ein steifes Feder-Dämpfer-Gelenk
  zwischen Segmenten (Riemendehnung als prismatische Feder). Näherung der ersten
  Mode, kein volles Modalverhalten.
\item \textbf{Echte Strukturvibration:} FEM / flexible Mehrkörpersimulation --
  SolidWorks (modal), MATLAB \textbf{Simscape Multibody}, Adams; Co-Simulation
  mit ROS möglich.
\item \textbf{RViz:} reine Visualisierung -- Kräfte als Marker, wenn ein anderer
  Knoten sie berechnet.
\end{itemize}
\begin{infobox}[title=Empfohlene Aufteilung]
\textbf{Gazebo} für Regelung, Trajektorientracking, Momentenbudget und
Ball-Schläger-Kontakt. \textbf{FEM/Simscape} für die strukturelle Resonanz
($f_1$). Beides verbindet die Input-Shaping-Frequenz.
\end{infobox}

\subsection{CAD-Module und URDF}
Wie man eine URDF aus Links, Joints und einem \code{ros2\_control}-Block aufbaut,
steht in Projekt~0.A (Kap.~\ref{ch:projekt0a}) und im CAD$\to$URDF-Kapitel
(Kap.~\ref{ch:cad_urdf}). Hier das Gantry-Spezifische.
\subsubsection{Notwendige CAD-Module mit Designanforderungen}
\begin{longtable}{@{}p{3.0cm}p{4.4cm}p{7.2cm}@{}}
\toprule
\textbf{Modul} & \textbf{Funktion} & \textbf{Designanforderung} \\
\midrule
Rahmen / Basis         & trägt X-Achse + X-Motor, ortsfest & sehr steif (Alu-Profil), trianguliert; URDF-Wurzel \\
X-Achse                & horizontale Verfahrachse & Hohlstahl-Stäbe + GT2-Riemen; Motor rahmenfest; $>1{,}6\,$m \\
X-Schlitten            & trägt + verfährt die Z-Achse & leicht, biegesteif; Stababstand $=$ Z-Hub \\
Z-Achse                & vertikale Verfahrachse & gleiches Prinzip; Hub $=$ X-Stababstand \\
Z-Schlitten            & trägt Y/Arm & geringe Masse (bewegte Masse!) \\
Y-Achse (optional)     & Tiefenzustellung & nur falls nötig; sonst starrer Ausleger \\
Armbasis               & Anbindung Arm an Schlitten & steife Einspannung; kurzer Hebel $L$ \\
Armglieder (2--3 DOF)  & Schwung + Schlägerwinkel & leicht (Alu/CFK); Motoren basisnah \\
Schlägereinspannung    & hält den Schläger & starr, definierte Schlägerebene \\
Endeffektor-Frame      & Schlagpunkt + Normale & TF-Frame: Schlägerflächenmitte + Normale \\
\bottomrule
\end{longtable}
\subsubsection{Kinematische Kette}
\begin{center}
\begin{tikzpicture}[node distance=3.5mm and 7mm,every node/.style={font=\scriptsize}]
\node[node] (w)  {world};
\node[node,right=of w]  (b)  {base\_link};
\node[node,right=of b]  (xc) {x\_carriage};
\node[node,right=of xc] (zc) {z\_carriage};
\node[node,right=of zc] (ab) {arm\_base};
\node[node,right=of ab] (l1) {link1};
\node[node,below=5mm of l1] (l2) {link2};
\node[node,left=of l2]  (pad) {paddle};
\draw[flow] (w) --node[above,font=\tiny]{fixed}(b);
\draw[flow] (b) --node[above,font=\tiny,align=center]{joint\_x\\prism.}(xc);
\draw[flow] (xc)--node[above,font=\tiny,align=center]{joint\_z\\prism.}(zc);
\draw[flow] (zc)--node[above,font=\tiny]{[joint\_y]}(ab);
\draw[flow] (ab)--node[above,font=\tiny,align=center]{arm1\\rev.}(l1);
\draw[flow] (l1)--node[right,font=\tiny]{arm2 rev.}(l2);
\draw[flow] (l2)--node[above,font=\tiny]{wrist rev.}(pad);
\end{tikzpicture}
\captionof{figure}{Kinematische Kette: drei (optional vier) Linear-/Drehübergänge
bis zur Armbasis, dann der Schlagarm bis zum Schläger.}
\end{center}
\subsubsection{URDF/Xacro-Gerüst}
\begin{lstlisting}[style=xml,caption={Gantry + Schlagarm als URDF (gekürzt)}]
<link name="base_link"> <!-- Wurzel: Rahmen fest mit der Welt --> </link>

<joint name="joint_x" type="prismatic">          <!-- X ueber die Plattenbreite -->
  <parent link="base_link"/> <child link="x_carriage"/>
  <origin xyz="0 0 1.0"/> <axis xyz="1 0 0"/>
  <limit lower="0" upper="1.5" effort="60" velocity="3.0"/>
</joint>
<link name="x_carriage"> ... </link>

<joint name="joint_z" type="prismatic">          <!-- Z vertikal, Hub = X-Stababstand -->
  <parent link="x_carriage"/> <child link="z_carriage"/>
  <axis xyz="0 0 1"/>
  <limit lower="-0.4" upper="0.0" effort="50" velocity="2.5"/>
</joint>
<link name="z_carriage"> ... </link>

<joint name="joint_arm1" type="revolute">        <!-- Schlagarm (arm2, wrist analog) -->
  <parent link="arm_base"/> <child link="link1"/>
  <axis xyz="0 1 0"/>
  <limit lower="-2.0" upper="2.0" effort="8" velocity="12"/>
</joint>
<link name="link1"> ... </link>

<joint name="joint_paddle" type="fixed">         <!-- Schlaeger-Endeffektor -->
  <parent link="link2"/> <child link="paddle"/>
  <origin xyz="0 0 0.12"/>
</joint>
<link name="paddle"> ... </link>
\end{lstlisting}
\begin{infobox}[title=Riemenübersetzung der Linearachsen]
Eine prismatische Riemenachse wandelt Motordrehung in Linearweg: \emph{Weg pro
Umdrehung} $=$ Zähnezahl $\times$ Riementeilung. Ein 20-Zahn-GT2-Rad
($2\,$mm Teilung) ergibt $40\,$mm/Umdrehung; mit $200\times16$ Microsteps folgt
$40/3200=0{,}0125\,$mm/Microstep. Diese Umrechnung Motor$\leftrightarrow$mm
gehört ins Hardware-Interface; in der URDF stehen am Joint die \emph{linearen}
Limits.
\end{infobox}

\subsection{Ablauf- und Parallellogik}
Wahrnehmung, Schätzung und Stellgrößen laufen \textbf{nebenläufig} mit sehr
unterschiedlichen Raten -- als separate ROS\,2-Knoten (teils auf verschiedenen
Rechnern), die nur über Topics gekoppelt sind und sich nicht gegenseitig
blockieren (Producer/Consumer).
\begin{longtable}{@{}p{3.4cm}p{2.8cm}p{8.0cm}@{}}
\toprule
\textbf{Knoten} & \textbf{Rate} & \textbf{Aufgabe} \\
\midrule
\code{ball\_detector} & Kamera-FPS (30--120\,Hz) & Stereo-Erkennung $\rightarrow$ 3D-Punkt (Pi\,4) \\
\code{ball\_tracker} (EKF) & pro Messung & Zustand updaten, Flug prädizieren \\
\code{strike\_planner} & pro EKF-Update & Auftreffpunkt + Zeit + Schlagbahn \\
\code{robot\_controller} & $\approx$\,100\,Hz & Trajektorie interpolieren, Stellgrößen \\
ESP32 (micro-ROS) & Echtzeit (kHz) & STEP/DIR-Erzeugung, Endstops \\
\bottomrule
\end{longtable}
\begin{infobox}[title={Frühe Reaktion statt Warten auf die perfekte Schätzung}]
Der \code{strike\_planner} veröffentlicht bereits nach den \emph{ersten}
EKF-Schätzungen eine vorläufige \code{/robot/target\_pose}; der Roboter fährt
sofort grob in Richtung Abfangzone. Jede neue Ballmessung verfeinert die
Schätzung -- die \code{target\_pose} wird laufend aktualisiert, und der Regler
folgt der bewegten Sollgröße. So wird die knappe Flugzeit voll genutzt, statt auf
eine endgültige Schätzung zu warten.
\end{infobox}

% =============================================================
\section{Vorgehensweise}
Wir bauen die \textbf{Wahrnehmung zuerst vollständig in Simulation} und erst
danach die Robotermechanik und -regelung. So lässt sich das Balltracking und die
Auftreffpunkt-Schätzung isoliert verifizieren (direkt auf Projekt~0.B
aufbauend), bevor die Mechanik dazukommt.

\subsection{Phase 1 -- Kamera und Balltracking in Simulation}

\subsubsection{Tischtennisplatte als CAD/SDF-Modell}
Zuerst wird die Tischtennisplatte konstruiert (CAD $\to$ Mesh/Primitive) und als
\emph{statisches} Modell in die Gazebo-Welt eingefügt: Platte
$2{,}74\times1{,}525\times0{,}76\,$m mit Netz. Sie dient als Referenzgeometrie
und definiert die \textbf{Zielebene} (siehe unten). Die Stereokamera-Welt aus
Projekt~0.B (Kap.~\ref{ch:projekt0b}) wird um dieses Modell ergänzt.

\subsubsection{Gegnerischen Ballschlag parametrisch simulieren}
Ein Knoten (oder Gazebo-Plugin) startet den Ball mit einer
\emph{parametrierbaren} Anfangsbedingung -- so sind Schläge reproduzierbar:
\begin{lstlisting}[style=yaml,caption={Parameter eines simulierten Gegnerschlags}]
start_position:  [x0, y0, z0]      # Abschlagort auf der Gegnerseite
speed:           v0                 # Betrag der Abschlaggeschwindigkeit [m/s]
azimuth:         phi                # horizontaler Abschlagwinkel [rad]
elevation:       theta              # vertikaler Abschlagwinkel [rad]
# optional (spaeter): spin_axis, spin_rate
\end{lstlisting}
\noindent Daraus folgt die Anfangsgeschwindigkeit
$\mathbf{v}_0=v_0\,(\cos\theta\cos\phi,\ \cos\theta\sin\phi,\ \sin\theta)$, die
dem Ball als Twist gesetzt wird; die Schwerkraft übernimmt Gazebo.

\subsubsection{Stereokamera-Tracking des Balls}
Die Wahrnehmungskette aus 0.B (\code{ball\_detector} $\rightarrow$
\code{/ball/point}) trackt den fliegenden Ball und liefert verrauschte
3D-Positionen im Kamerasystem.

\subsubsection{Auftreffpunkt über schiefen Wurf und EKF}
Ohne Luftwiderstand ist der Flug ein \textbf{schiefer Wurf} mit bekannter
Erdbeschleunigung $\mathbf{g}=(0,0,-g)$:
\begin{equation}
\mathbf{r}(t) = \mathbf{r}_0 + \mathbf{v}_0\,t + \tfrac{1}{2}\,\mathbf{g}\,t^2 .
\end{equation}
Der EKF-Zustand ist $\mathbf{x}=[\mathbf{r},\mathbf{v}]^\top$ (6D), Messung sind
die 3D-Stereopositionen, das Prozessmodell ist die konstante Beschleunigung
$\mathbf{g}$. Die Ballmasse ist bekannt -- sie wird erst relevant, sobald
Luftwiderstand/Spin das Modell \emph{nichtlinear} machen (dann ist der EKF
zwingend, beim reinen Wurf genügte ein linearer KF). Mit jeder neuen Messung
werden die Wurfparameter (effektiv $\mathbf{r}_0,\mathbf{v}_0$ bzw.\ der
aktuelle Zustand) aktualisiert; reine Prädiktionsschritte extrapolieren den Flug.

\paragraph{Schätzung nur bis zur Tischmitte -- danach unveränderbar.} Der EKF
korrigiert die Bahn nur, solange der Ball auf der \emph{Gegnerseite} ist
(bis zur Tischmitte, $y=0$). Überschreitet die Schätzung die Mitte, wird sie
\textbf{eingefroren} -- Abfangpunkt und -zeit gelten ab dann als fest. Drei
Gründe: (1)~Der Roboter beginnt eine \emph{frühe} Abfangbewegung und verfeinert
sie bis zur Mitte, fährt danach aber \emph{ruhig} zu Ende -- statt kurz vor dem
Auftreffpunkt noch schnelle, mechanisch kaum umsetzbare Korrekturen zu machen.
(2)~In der zeitkritischen Endphase fällt die EKF-Rechenlast weg. (3)~Bis zur
Mitte ist die Schätzung ohnehin konvergiert. Ein neuer Schlag (Ball wieder weit
auf der Gegnerseite) setzt das Filter zurück.

\paragraph{Aufprall auf der Platte (Restitution).} Trifft der Ball die
Tischfläche, ist das ein \emph{Stoß} statt einer glatten Parabel. Ein realer
Tischtennisball verliert dabei wenig Energie: aus der ITTF-Prüfung (Fall aus
$30{,}5\,$cm, Rücksprung $24$--$26\,$cm) folgt ein Restitutionskoeffizient
$e=\sqrt{h_{\text{zurück}}/h_{\text{fall}}}\approx0{,}9$ -- \emph{nahezu}
elastisch, aber nicht ideal ($e=1$). Im Flugmodell wird der Aufprall als
\textbf{stückweise Reflexion} eingebaut: schneidet die prädizierte Bahn die
Tischebene ($z=h_{\text{Tisch}}$) innerhalb der Platte, kehrt sich die
Normalkomponente um und wird gedämpft,
\begin{equation}
v_z \;\rightarrow\; -e\,v_z,\qquad (v_x,v_y\ \text{bei Reibungsfreiheit unverändert}),
\end{equation}
danach läuft die Parabel weiter. \textbf{Reibung} beim Aufprall wirkt nur
tangential und damit fast ausschließlich auf den \emph{Spin}; auf die
translatorische Bahn ist sie ohne Spin vernachlässigbar. Ein ideal elastischer
Stoß ($e=1$) ist eine brauchbare erste Näherung, $e\approx0{,}9$ ist genauer und
verschiebt Rücksprunghöhe und damit Abfangzeit spürbar. In der Simulation liefert
Gazebo dieses Verhalten über den \code{restitution\_coefficient} der
Kollisionsfläche (Schritt~2).

\subsubsection{Zielebene, Plattenrand und Fehlschuss-Erkennung}
Zwei Ebenen sind relevant. \textbf{(1)~Die Tischfläche}
($z=h_{\text{Tisch}}=0{,}76\,$m, horizontal): Der Schnittpunkt der prädizierten
Bahn mit ihr ist der \emph{Aufprallpunkt} auf der Platte. Liegt er \emph{innerhalb}
des Plattenrands auf der Roboterhälfte, ist der Schlag gültig (dort findet der
Restitutions-Bounce statt); liegt er außerhalb, geht der Ball ins Aus.
\textbf{(2)~Die Abfangebene} (vertikal, z.\,B.\ $y=y_R$ an der hinteren
Plattenkante der Roboterseite): ihr Durchstoßpunkt \emph{nach} dem Bounce ist die
eigentliche \textbf{Abfangposition} des Schlägers,
\begin{equation}
t^\ast=\frac{y_R-y_0}{v_{y0}},\qquad x^\ast=x(t^\ast),\quad z^\ast=z(t^\ast),
\end{equation}
also $(x^\ast,y_R,z^\ast)$ und Abfangzeit $t^\ast$ (mit dem gespiegelten
Geschwindigkeitsvektor nach dem Tischkontakt). Als \textbf{Fehlschuss} markiert
wird (Flag/Diagnostic): wenn der Ball das gemeinsame Sichtfeld beider Kameras
verlässt (keine Triangulation mehr -- vgl.\ FOV-Kegel, Kap.~\ref{ch:projekt0b}),
\emph{oder} der Aufprallpunkt außerhalb der Platte liegt, \emph{oder} die
Abfangposition außerhalb des Arbeitsraums, \emph{oder} ein \textbf{Netztreffer}
vorhergesagt wird.

\paragraph{Netztreffer und Netzroller.} Schneidet die prädizierte Bahn die
\emph{Netzebene} ($y=0$) \emph{unterhalb} der Netzoberkante ($z<0{,}9125\,$m),
geht der Ball ins Netz -- der Wurf wird ebenfalls als \textbf{ungültig} erkannt.
Auch ein \emph{Netzroller} (der Ball berührt die Netzkante und fällt knapp
dahinter) gilt als nicht abfangbar: die dabei stark verlangsamte, abgelenkte Bahn
lässt dem Roboter weder Bewegungsraum noch Reaktionszeit, ihn regulär zu treffen.
Solche Bälle werden gar nicht erst angefahren.

\paragraph{Aufschlag- vs.\ Ballwechselregeln.} Bisher haben die Beispielwürfe nur
einen \emph{Ballwechsel} nachgebildet (der Ball kommt über das Netz und trifft
zuerst die Roboterseite). Ein regelkonformer \textbf{Aufschlag} sieht anders aus:
Der Ball muss \emph{zuerst genau einmal auf der Aufschlagseite} ($y>0$, beim
Gegenspieler) aufkommen und \emph{danach} auf der Roboterseite ($y<0$). Erst dann
ist der Aufschlag gültig. \textbf{Der Roboter schlägt nie selbst auf} -- der
Aufschlag kommt immer vom Spieler; der Roboter ist hier ausschließlich
Rückschläger. Daraus folgen die Fehlschuss-Fälle des Aufschlags: erster Aufprall
\emph{nicht} auf der eigenen Seite, ein \emph{zweiter} Aufprall erneut auf der
Aufschlagseite, oder ein Netztreffer. Nach einem gültigen Aufschlag wird auf die
\textbf{Ballwechselregeln} umgeschaltet: der ankommende Ball muss genau einmal auf
der Roboterseite aufkommen (ein zweiter Aufprall ohne Rückschlag ist ein
Doppelaufprall). Dieselbe kleine Zustandsmaschine, die schon zwischen REAKTIV und
PAUSE trennt, verwaltet damit auch die Phase \emph{Aufschlag} bzw.\
\emph{Ballwechsel}.

\paragraph{Sichtbare Begrenzungswand als Güte-Anzeige.} Um die Qualität von EKF
und Flugmodell \emph{direkt} beurteilen zu können, wird am Plattenrand der
Roboterseite (an der Abfangebene $y=y_R$) eine \textbf{durchsichtige
Begrenzungswand} eingezogen -- rein visuell, der Ball fliegt hindurch. Auf ihr
wird der vom EKF geschätzte \emph{Auftreffpunkt} als Marker eingeblendet, der sich
mit jeder Messung aktualisiert. So sieht man unmittelbar, \emph{wo} das Modell den
Durchstoßpunkt erwartet und \emph{wo} der reale Ball die Wand tatsächlich
durchbricht -- die Abweichung ist die Modellgüte. Ein \textbf{Fehlschuss} wird an
derselben Stelle sichtbar gemacht (die Wand färbt sich rot) und beim nächsten
Ballwechsel wieder zurückgesetzt.

\paragraph{Abbruch und Rückkehr in die Initialstellung.} Wird ein Wurf -- noch
\emph{während} des Anfahrens -- als ungültig erkannt (Aus, Netz, außer
Reichweite), muss die laufende Roboterbewegung \textbf{abgebrochen} werden: das
aktive Goal wird gecancelt, der Roboter stoppt die Abfangbewegung und fährt in
seine \textbf{Initialstellung} (Home) zurück -- bereit für den nächsten
Ballwechsel. Genau dafür ist der \textbf{Action}-Mechanismus gemacht
(Projekt~0.A, Kap.~\ref{ch:projekt0a}): ein Goal ist jederzeit \emph{cancelbar},
und der Action-Server reagiert mit einem definierten Abbruch -- in Simulation und
Realität identisch.

\paragraph{Betriebsmodi: reaktiv vs.\ Pause.} Der Roboter darf \emph{nicht
immer} auf erkannte Bälle reagieren. Eine kleine \textbf{Zustandsmaschine} trennt:
\begin{itemize}
\item \textbf{REAKTIV} -- ein Ballwechsel läuft: die Tracker-Schätzungen werden in
  Roboter-Goals umgesetzt (frühe Abfangbewegung, laufende Verfeinerung).
\item \textbf{PAUSE} -- zwischen den Ballwechseln (Ball holen, Aufschlag
  vorbereiten): der Tracker darf weiterlaufen, löst aber \emph{keine}
  Roboterbewegung aus. So werden \textbf{Fehlreaktionen} ausgeschlossen, etwa wenn
  jemand den Ball aufhebt oder über den Tisch greift.
\end{itemize}
Diese Moduslogik gehört bereits in den \emph{Entwurf}. Da die Roboterbefehle
ohnehin über \textbf{Actions und Goals} laufen, ist das passende Konstrukt schon
vorhanden: im PAUSE-Modus werden schlicht keine Goals gesendet bzw.\ angenommen,
im REAKTIV-Modus sendet der \code{strike\_planner} Goals und
aktualisiert bzw.\ cancelt sie. Derselbe Mechanismus trägt von der Simulation in
die reale Anlage.

\subsubsection{Ausblick: Luftwiderstand}
Beim leichten Tischtennisball (Masse $m\approx2{,}7\,$g, Durchmesser $40\,$mm)
ist der Luftwiderstand \emph{nicht} vernachlässigbar -- das kleine Verhältnis von
Masse zu Querschnittsfläche führt zu spürbarer Verzögerung. Die
Widerstandskraft wirkt entgegen der Geschwindigkeit und wächst quadratisch:
\begin{equation}
\mathbf{F}_D = -\tfrac{1}{2}\,\rho\,C_D\,A\,\lVert\mathbf{v}\rVert\,\mathbf{v},
\qquad A=\pi r^2,
\end{equation}
mit Luftdichte $\rho\approx1{,}2\,$kg/m\textsuperscript{3}, Widerstandsbeiwert
$C_D\approx0{,}4$ (Kugel) und Querschnittsfläche $A$. Die Bewegungsgleichung
lautet damit
\begin{equation}
\dot{\mathbf{v}} = \mathbf{g} - k\,\lVert\mathbf{v}\rVert\,\mathbf{v},
\qquad k=\frac{\rho\,C_D\,A}{2m},
\end{equation}
und ist \textbf{nichtlinear} in $\mathbf{v}$ -- genau deshalb wird hier der
\textbf{EKF} (statt eines linearen KF) gebraucht: Das Prozessmodell wird pro
Schritt linearisiert (Jacobi-Matrix der Beschleunigung nach dem Zustand). Die
bekannte Ballmasse $m$ geht direkt in $k$ ein. Wirkung: Die reale Flugbahn ist
\emph{flacher und kürzer} als die reine Wurfparabel -- der Auftreffpunkt
verschiebt sich, weshalb der Term für eine genaue Abfangzeit-Schätzung relevant
ist. Der Schritt vom reinen Wurf zum Modell mit Luftwiderstand ist die natürliche
erste Verfeinerung; der Spin (folgender Abschnitt) baut darauf auf.

\subsubsection{Ausblick: Ballspin (Magnus)}
Das Modell wird so aufgebaut, dass es um \textbf{Spin} erweiterbar ist: Die
Magnus-Kraft $\mathbf{F}_M\propto \boldsymbol\omega\times\mathbf{v}$ ergänzt das
Prozessmodell (nichtlinear $\rightarrow$ EKF/UKF), der Zustand wird um die
Drehrate $\boldsymbol\omega$ erweitert. Messbar wird der Spin, indem
\textbf{Markierungen auf dem Ball} getrackt werden, um Rotationsrichtung und
-geschwindigkeit zu schätzen -- eine spätere Ausbaustufe.

\subsection{Phase 2 -- Roboteranbindung}

\subsubsection{Frühe Abfangbewegung (iterative Verfeinerung)}
Gemäß der Ablauflogik fährt der Roboter bereits nach den ersten Schätzungen los
(grobe \code{target\_pose}) und verfeinert das Ziel mit jeder neuen Messung. Das
minimiert die effektive Reaktionszeit innerhalb der kurzen Ballflugdauer.

\subsubsection{Stellgrößen und Motorsteuerung}
Die geschätzte \code{target\_pose} (Abfangpunkt + Schlägerorientierung) wird per
inverser Kinematik in Achswerte umgesetzt: $x,z$ (und optional $y$) aus der
kartesischen Zielposition, die Armgelenke für den Schlägerwinkel. Der
\code{joint\_trajectory\_controller} interpoliert, das Hardware-Interface
(micro-ROS $\rightarrow$ ESP32 $\rightarrow$ TMC2209) erzeugt die Schritte --
derselbe \code{ros2\_control}-Stack wie in Projekt~0.A und Projekt~1. Hier zählt
das \textbf{Timing}: in der Simulation trivial, real der kritische Faktor.

% =============================================================
\section{Umsetzung}
Hier folgen die konkreten Implementierungsschritte -- jeder für sich lauffähig,
bevor der nächste daraufgesetzt wird. Wie in den Einstiegsprojekten lebt der Code
in einem eigenen Paket \code{p2\_tabletennis} im gemeinsamen Workspace
\code{\textasciitilde/ros2\_ws}.

\subsection{Schritt 1 -- Tischtennisplatte als STL in Gazebo einbinden}
Das CAD-Modell der Platte (aus SolidWorks o.\,Ä.\ als \code{.stl} exportiert)
wird als statisches Mesh-Modell in eine Gazebo-Welt eingebunden. Zwei Eigenheiten
von STL-Dateien sind dabei zu beachten:
\begin{itemize}
\item \textbf{Einheiten:} STL trägt keine Einheit. Der Export ist hier in
  \emph{Millimetern} (Bounding-Box $1525\times2740\times912$) -- Gazebo rechnet in
  \emph{Metern}, also \code{<scale>0.001 0.001 0.001</scale>}.
\item \textbf{Orientierung:} Im CAD zeigt $Y$ nach oben (die Platte sitzt auf der
  $y=0$-Ebene), in ROS/Gazebo zeigt $Z$ nach oben. Eine Modell-\code{<pose>} mit
  Roll $+90^\circ$ ($1{,}5708$) dreht $Y$-oben auf $Z$-oben.
\end{itemize}
\begin{lstlisting}[style=xml,caption={\code{worlds/tabletennis.sdf} -- die Platte als Mesh-Modell}]
<model name="table">
  <static>true</static>
  <pose>0.173 0 0 1.5708 0 0</pose>   <!-- x: Breite zentrieren; roll +90: Y->Z -->
  <link name="link">
    <visual name="v"><geometry><mesh>
      <uri>package://p2_tabletennis/meshes/Tischtennisplatte.stl</uri>
      <scale>0.001 0.001 0.001</scale>            <!-- mm -> m -->
    </mesh></geometry></visual>
    <collision name="c"><geometry><mesh>
      <uri>package://p2_tabletennis/meshes/Tischtennisplatte.stl</uri>
      <scale>0.001 0.001 0.001</scale>
    </mesh></geometry></collision>
  </link>
</model>
\end{lstlisting}

\begin{warnbox}[title={Gazebo löst \code{package://} nur über \code{GZ\_SIM\_RESOURCE\_PATH}}]
Gazebo kennt keine ROS-Pakete. Damit \code{package://p2\_tabletennis/\ldots}
auflöst, hängt die Launch-Datei den \emph{Eltern-Ordner} des Paket-share an den
Suchpfad -- Gazebo sucht dort den ersten URI-Teil als Unterordner. Deshalb
\textbf{immer über die Launch-Datei starten}, nicht mit nacktem \code{gz sim}.
\begin{lstlisting}[style=py,caption={\code{launch/table.launch.py} (Kern)}]
from launch.actions import AppendEnvironmentVariable
pkg_share = get_package_share_directory('p2_tabletennis')
AppendEnvironmentVariable('GZ_SIM_RESOURCE_PATH', os.path.dirname(pkg_share))
# ... IncludeLaunchDescription(gz_sim.launch.py, gz_args=[world, ' -r'])
\end{lstlisting}
\end{warnbox}

\begin{lstlisting}[style=bash,caption={Bauen, sourcen, Welt starten}]
cd ~/ros2_ws
colcon build --packages-select p2_tabletennis
source install/setup.bash
ros2 launch p2_tabletennis table.launch.py
\end{lstlisting}

\begin{figure}[htbp]\centering
  \includegraphics[width=0.85\textwidth]{project_2/tabletennis_plate.png}
  \caption{Die eingebundene Tischtennisplatte in Gazebo: korrekt in Metern
  skaliert (Spielfläche bei $z=0{,}76\,$m) und über Roll\,$+90^\circ$ aufrecht
  gestellt. Sie bildet die Referenzgeometrie und die spätere Zielebene.}
  \label{fig:p2-table}
\end{figure}

\begin{lstlisting}[style=bash,caption={\code{ros2\_ws}-Delta nach Schritt 1 (\# neu / \# angepasst)}]
src/p2_tabletennis/
|- meshes/Tischtennisplatte.stl   # neu (STL aus dem CAD, in mm)
|- worlds/tabletennis.sdf         # neu (Platte als Mesh-Modell + Boden + Licht)
|- launch/table.launch.py         # neu (Gazebo + GZ_SIM_RESOURCE_PATH)
`- setup.py                       # angepasst: data_files (worlds/launch/meshes)
\end{lstlisting}

\subsection{Schritt 2 -- Stereo-Rig und Ball in die Welt integrieren}
Die Wahrnehmung aus Projekt~0.B (Kap.~\ref{ch:projekt0b}) wird übernommen: das
Stereo-Rig und der Ball kommen in \code{tabletennis.sdf}, die Knoten
(\code{ball\_detector}, \code{ball\_tracker}, \code{marker\_viz},
\code{fov\_viz}) werden unverändert aus \code{p0b\_stereo} wiederverwendet.

\paragraph{Welt erweitern.} Stereo-Rig (Blick entlang $+y$ zur Gegnerseite,
Basisbreite $0{,}10\,$m) und der orange Ball (40\,mm, 2{,}7\,g) kommen hinzu.
Damit der Ball \textbf{nicht durch die Platte fällt}, bekommt die Spielfläche
eine \emph{Box-Kollision} (robuster als ein konkaves Mesh) mit Restitution:
\begin{lstlisting}[style=xml,caption={Spielflächen-Kollision + Ball (Auszug aus tabletennis.sdf)}]
<model name="table_top"><static>true</static>
  <pose>0 0 0.74 0 0 0</pose>                 <!-- Oberkante z=0.76 -->
  <link name="link"><collision name="c">
    <geometry><box><size>1.525 2.74 0.04</size></box></geometry>
    <surface>
      <bounce><restitution_coefficient>0.9</restitution_coefficient><threshold>0</threshold></bounce>
      <friction><ode><mu>0.4</mu><mu2>0.4</mu2></ode></friction>
    </surface>
  </collision></link>
</model>

<model name="ball"><pose>0 0.5 1.2 0 0 0</pose>
  <link name="link">
    <inertial><mass>0.0027</mass> ... </inertial>
    <collision name="c"><geometry><sphere><radius>0.02</radius></sphere></geometry>
      <surface><bounce><restitution_coefficient>0.9</restitution_coefficient>
        <threshold>0</threshold></bounce></surface></collision>
    <visual name="v"> ... orange ... </visual>
  </link>
</model>
\end{lstlisting}
\begin{infobox}[title={Aufprall: kein Durchfallen, fast elastischer Stoß}]
Eine \emph{Box}-Kollision der Spielfläche (statt des konkaven Tisch-Meshes)
verhindert Tunneling; mit dem kleinen Zeitschritt ($1\,$ms) bleibt der Ball
sicher auf dem Tisch (Test: Start $z=1{,}2$, nach Sekunden noch springend bei
$z\approx0{,}83$ -- nicht durchgefallen). Der \code{restitution\_coefficient}
$0{,}9$ entspricht dem realen Ball (nahezu elastisch, vgl.\ 26.2.1.4); Reibung
beim Aufprall ist ohne Spin vernachlässigbar.
\end{infobox}

\paragraph{Netz: Box-Kollision statt eigenem STL.} Das Tisch-Mesh enthält das
Netz nur \emph{visuell} -- ohne Kollision fliegt der Ball hindurch. Ein eigener
STL-Körper fürs Netz ist dafür \emph{nicht} nötig (und als dünner, konkaver
Collider sogar tunneling-anfällig); robuster ist -- wie bei der Spielfläche --
eine dünne \emph{Box}-Kollision am Netzort ($y=0$, Höhe $15{,}25\,$cm über der
Platte):
\begin{lstlisting}[style=xml,caption={Netz-Kollision (Auszug aus tabletennis.sdf)}]
<model name="net"><static>true</static>
  <pose>0 0 0.83625 0 0 0</pose>            <!-- 0.76 + 0.1525/2 -->
  <link name="link"><collision name="c">
    <geometry><box><size>1.525 0.01 0.1525</size></box></geometry>
    <surface>
      <bounce><restitution_coefficient>0.1</restitution_coefficient><threshold>0</threshold></bounce>
      <friction><ode><mu>0.8</mu><mu2>0.8</mu2></ode></friction>
    </surface>
  </collision></link>
</model>
\end{lstlisting}
\noindent Die \textbf{niedrige Restitution} ($0{,}1$) und \textbf{hohe Reibung}
nähern das weiche, energiezehrende Verhalten des Netzes an: der Ball prallt nicht
wie von einer Wand ab, sondern verliert Energie und fällt. Eine genauere
\emph{Netzphysik} (nachgiebiges Gewebe) bräuchte ein Soft-Body- bzw.\
Feder-Masse-Netz -- in Gazebo aufwändig; für die Spiel-Logik genügt meist ein
\emph{Netztreffer-Flag} (Let/Fehler) zusätzlich zur Fehlschuss-Erkennung
(26.2.1.5).

\paragraph{Sichtkegel.} Wie in Projekt~0.B (Kap.~\ref{ch:projekt0b},
Sichtfeld-Visualisierung) zeichnet \code{fov\_viz} die Sichtkegel beider Kameras
als \code{MarkerArray} auf \code{/stereo/fov\_markers}; ihr Überlappungsbereich
ist die Zone, in der trianguliert werden kann.

\paragraph{Pipeline per Launch.} Eine Launch-Datei bündelt -- mit allen Lehren
aus 0.A/0.B -- Gazebo, \code{/clock}-Brücke, Bildbrücke, statische Kamera-TF und
die vier Knoten, alle mit \code{use\_sim\_time}:
\begin{lstlisting}[style=py,caption={\code{launch/tabletennis.launch.py} (gekürzt)}]
detector = Node(package='p0b_stereo', executable='ball_detector', parameters=[SIM])
tracker  = Node(package='p0b_stereo', executable='ball_tracker',  parameters=[SIM])
viz      = Node(package='p0b_stereo', executable='marker_viz',    parameters=[SIM])
fov      = Node(package='p0b_stereo', executable='fov_viz',       parameters=[SIM])
static_tf = Node(package='tf2_ros', executable='static_transform_publisher',
    arguments=['--x','0','--y','-1.9','--z','1.1','--roll','-1.5708',
               '--frame-id','world','--child-frame-id','stereo_left_optical'])
return LaunchDescription([resource, gz, clock_bridge, image_bridge,
                          static_tf, detector, tracker, viz, fov, rviz])
\end{lstlisting}
\begin{infobox}[title={Checkliste 0.A/0.B -- in dieser Launch beachtet}]
\code{/clock} gebrückt \emph{und} \code{use\_sim\_time:=true} auf allen Knoten
(sonst ``No transform''); statische TF \code{world}$\rightarrow$%
\code{stereo\_left\_optical}, damit der Kamerarahmen existiert; \code{ogre2}/EGL
fürs Kamerarendering; \code{ApproximateTimeSynchronizer} für L/R-Bilder; Ball
orange ($\Rightarrow$ HSV-Bereich aus 0.B passt); \code{package://} über
\code{GZ\_SIM\_RESOURCE\_PATH}.
\end{infobox}

\paragraph{Terminals und Befehle.}
\begin{lstlisting}[style=bash,caption={Terminal 1 -- bauen + alles starten (Gazebo + RViz + Pipeline)}]
cd ~/ros2_ws
colcon build --packages-select p2_tabletennis
source install/setup.bash
ros2 launch p2_tabletennis tabletennis.launch.py
\end{lstlisting}
\begin{lstlisting}[style=bash,caption={Terminal 2 -- Ballposition anzeigen}]
source ~/ros2_ws/install/setup.bash
ros2 topic echo /ball/point        # gemessene 3D-Position (Triangulation)
ros2 topic echo /ball/estimate     # Kalman-Schaetzung
\end{lstlisting}
\noindent In Gazebo erscheinen Platte, Rig und der (fallende/springende) Ball;
in RViz die Sichtkegel, die Marker (gemessen/geschätzt/Vorhersage) und beide
Kamerabilder. Der \emph{parametrische Gegnerschlag} (Anfangsgeschwindigkeit statt
freier Fall) folgt später; bis dahin lässt sich der Ball in Gazebo mit der Maus
bewegen.
\begin{lstlisting}[style=bash,caption={\code{ros2\_ws}-Delta nach Schritt 2 (\# neu / \# angepasst)}]
src/p2_tabletennis/
|- worlds/tabletennis.sdf       # angepasst: Stereo-Rig + Ball + Box-Kollision (Platte + Netz)
|- launch/tabletennis.launch.py # neu (Welt + Bruecken + Pipeline + RViz)
|- config/tabletennis.rviz      # neu (Marker, FOV, Kamerabilder; Fixed Frame world)
`- package.xml                  # angepasst: exec_depend p0b_stereo u.a.
\end{lstlisting}

\subsection{Schritt 3 -- Balltracking: fliegender Ball}
Damit die Pipeline einen echten Flug verfolgt, muss der Ball über den Tisch
fliegen -- der Gegnerschlag aus 26.2.1.2. Gazebo bietet \emph{keine} API, einem
freien Körper eine Anfangsgeschwindigkeit zu geben (nur \code{set\_pose}). Wir
bewegen den Ball daher \textbf{kinematisch} entlang einer parametrischen
Wurfparabel: der Knoten \code{opponent\_serve} integriert Position und
Geschwindigkeit (mit Tisch-Bounce, vgl.\ 26.2.1.4) und setzt die Ballpose
$\sim$50x/s über den gebrückten \code{SetEntityPose}-Service. Vorteil:
reproduzierbar und mit bekannter Ground Truth -- ideal, um Erkennung und EKF
gegenzuprüfen. (\code{self.p(n)} unten ist Kürzel für
\code{self.get\_parameter(n).value}.)
\begin{lstlisting}[style=py,caption={\code{opponent\_serve} -- Kern (Wurf + Bounce)}]
def start_serve(self):                      # v0 aus speed/azimuth/elevation
    s, sp = self.p('start'), self.p('speed')
    az, el = self.p('azimuth'), self.p('elevation')
    self.pos = [float(s[0]), float(s[1]), float(s[2])]
    self.vel = [sp*cos(el)*cos(az), sp*cos(el)*sin(az), sp*sin(el)]
    self.last_t = None

def tick(self):
    if self.pos is None: return
    now = self.get_clock().now().nanoseconds * 1e-9    # dt aus (Sim-)Zeit!
    if self.last_t is None: self.last_t = now; return
    dt = now - self.last_t; self.last_t = now
    self.vel[2] -= self.p('g') * dt
    for i in range(3): self.pos[i] += self.vel[i] * dt
    if self.pos[2] <= TZ and self.vel[2] < 0 and on_table(self.pos):
        self.pos[2] = TZ; self.vel[2] *= -self.p('restitution')   # Bounce
    self.set_ball(self.pos)                 # SetEntityPose-Service (async)
\end{lstlisting}

\begin{infobox}[title={Lehre: dt aus der Uhr -- und \code{use\_sim\_time} braucht /clock}]
Mit \emph{festem} \code{dt} pro Tick flog der Ball zu schnell, weil der Timer
unter \code{use\_sim\_time} beim Sim-Zeit-Aufholen mehrfach feuert; korrekt ist
\code{dt} aus der \emph{tatsächlichen} Uhr-Differenz. Und: läuft
\code{opponent\_serve} mit \code{use\_sim\_time}, aber \emph{ohne} gebrückte
\code{/clock}, feuert sein Timer nie -- der Ball bleibt stehen (genau der
Checklistenpunkt aus 24.6). Die \code{/clock}-Brücke ist in der Launch deshalb
zwingend.
\end{infobox}

\noindent Den fliegenden Ball verarbeitet die \textbf{unveränderte Pipeline aus
Schritt 2}: \code{ball\_detector} (HSV + sparse Triangulation, 0.B) $\rightarrow$
\code{/ball/point}; \code{ball\_tracker} (Kalman) $\rightarrow$
\code{/ball/estimate} und \code{/ball/prediction}; \code{marker\_viz} zeigt alles
in RViz. Verlässt der Ball das gemeinsame Sichtfeld, stoppt \code{/ball/point} --
die Marker frieren auf der letzten Position ein (Fehlschuss-Indiz, vgl.\
26.2.1.5).

\begin{infobox}[title={Checkliste 23.5/24.6 -- in Schritt 3 beachtet}]
\code{/clock} + \code{use\_sim\_time} (sonst kein Servetimer); Pose-Felder mit
\code{float(\ldots)} gecastet (NumPy-Skalare in ROS-Feldern); \emph{kein}
Custom-Interface -- Auslösung über \code{std\_srvs/Trigger}; der Callback heißt
\code{on\_serve}, nicht \code{handle} (rclpy-Kollision); HSV-Bereich,
L/R-Synchronisation und Disparität unverändert aus 0.B.
\end{infobox}

\paragraph{Terminals und Befehle.}
\begin{lstlisting}[style=bash,caption={Terminal 1 -- alles starten (Welt + Pipeline + Gegnerschlag + RViz)}]
cd ~/ros2_ws
colcon build --packages-select p2_tabletennis
source install/setup.bash
ros2 launch p2_tabletennis tabletennis.launch.py
\end{lstlisting}
\begin{lstlisting}[style=bash,caption={Terminal 2 -- Messung und Schaetzung mitlesen}]
source ~/ros2_ws/install/setup.bash
ros2 topic echo /ball/point        # gemessene 3D-Position (Triangulation)
ros2 topic echo /ball/estimate     # Kalman-Schaetzung
\end{lstlisting}
\begin{lstlisting}[style=bash,caption={Terminal 3 (optional) -- Schlag manuell ausloesen}]
source ~/ros2_ws/install/setup.bash
ros2 service call /opponent/serve std_srvs/srv/Trigger "{}"
\end{lstlisting}
\noindent \code{opponent\_serve} schlägt standardmäßig alle $5\,$s automatisch
auf (Parameter \code{auto\_period}); \code{speed}, \code{azimuth},
\code{elevation} und \code{start} lassen sich anpassen, damit der Ball auf der
Roboterhälfte aufkommt und im Sichtfeld bleibt.
\begin{lstlisting}[style=bash,caption={\code{ros2\_ws}-Delta nach Schritt 3 (\# neu / \# angepasst)}]
src/p2_tabletennis/
|- p2_tabletennis/opponent_serve.py  # neu (parametrischer Gegnerschlag)
|- config/set_pose_bridge.yaml       # neu (Service-Bruecke set_pose)
|- launch/tabletennis.launch.py      # angepasst: set_pose-Bruecke + opponent_serve
|- setup.py                          # angepasst: entry_point opponent_serve
`- package.xml                       # angepasst: ros_gz_interfaces, std_srvs ...
\end{lstlisting}

\subsection{Schritt 4 -- Flugschätzung mit EKF und Luftwiderstand}
Der lineare 0.B-Tracker wird durch ein \textbf{EKF} ersetzt
(\code{flight\_ekf}), das den nichtlinearen Flug aus 26.2.1.4/26.2.1.6 schätzt:
Zustand $\mathbf{x}=[\mathbf{r},\mathbf{v}]^\top$ (6D), Prozessmodell
\begin{equation}
\dot{\mathbf r}=\mathbf v,\qquad
\dot{\mathbf v}=\mathbf g-k\,\lVert\mathbf v\rVert\,\mathbf v,\qquad
k=\frac{\rho\,C_D\,A}{2m}\approx0{,}131\,\mathrm{m^{-1}}.
\end{equation}
Wegen des $\lVert\mathbf v\rVert\mathbf v$-Terms ist das nichtlinear -- der EKF
linearisiert pro Schritt mit der Jacobi-Matrix; für den Geschwindigkeitsblock gilt
\begin{equation}
\frac{\partial \dot{\mathbf v}}{\partial \mathbf v}
= -k\Big(\lVert\mathbf v\rVert\,I+\frac{\mathbf v\mathbf v^\top}{\lVert\mathbf v\rVert}\Big).
\end{equation}
\begin{lstlisting}[style=py,caption={\code{flight\_ekf} -- EKF-Kern (Position gemessen, H=[I\ 0])}]
def ekf_step(self, z, dt):
    F = self._F(self.x, dt)               # Jacobi der Dynamik
    self.x = self._f(self.x, dt)          # Praediktion (nichtlinear, mit Drag)
    self.P = F @ self.P @ F.T + self.Q
    S = self.P[:3, :3] + self.R           # Korrektur, H = [I 0]
    K = self.P[:, :3] @ np.linalg.inv(S)
    self.x = self.x + K @ (z - self.x[:3])
    self.P = (np.eye(6) - K @ H) @ self.P
\end{lstlisting}

\begin{warnbox}[title={Physik im Welt-Frame -- Messungen per TF transformieren}]
Gravitation, Tischebene ($z=0{,}76$) und Abfangebene leben im \textbf{Welt-Frame},
die Messungen (\code{/ball/point}) kommen aber im \emph{optischen Kameraframe}.
Der EKF transformiert sie daher per TF nach \code{world}
(\code{tf\_buffer.transform(msg, 'world')}) und rechnet/publiziert in \code{world}.
In 0.B war das egal (konstante Geschwindigkeit, kein Schwerefeld) -- hier ist es
zwingend.
\end{warnbox}

\paragraph{Auftreffpunkt vorhersagen.} Nach jedem Update wird der Zustand
\emph{ohne} Messung vorwärts integriert (mit Tisch-Bounce) bis zur Abfangebene
$y=y_R$; das ergibt Abfangpunkt und -zeit. Netztreffer und Aus werden dabei wie
in 26.2.1.5 erkannt.
\begin{lstlisting}[style=py,caption={\code{predict\_impact} -- Bounce, Netz, Abfangebene}]
def predict_impact(self):
    x = self.x.copy(); pts = [x[:3].copy()]
    for _ in range(STEPS):
        yprev = x[1]; x = self._f(x, dt)
        if x[2] <= TZ and x[5] < 0 and on_table(x):   # Tisch-Bounce
            x[2] = TZ; x[5] *= -E
        if yprev > 0 >= x[1] and x[2] < NET_TOP:       # Netztreffer -> ungueltig
            return pts, None
        if x[1] <= yR:                                 # Abfangebene erreicht
            return pts, x[:3].copy()
    return pts, None
\end{lstlisting}
\noindent Veröffentlicht werden \code{/ball/estimate} und \code{/ball/prediction}
(wie 0.B, sodass \code{marker\_viz} weiterläuft) sowie \code{/ball/intercept} und
ein Marker (\code{/ball/intercept\_marker}, gelber Zylinder in RViz).

\paragraph{Einfrieren an der Tischmitte (Umsetzung).} Gemäß 26.2.1.4 korrigiert
der EKF nur bis zur Mitte ($y=0$); danach steht die Schätzung fest und es wird
\emph{nicht weitergerechnet} (Performance, ruhige Endanfahrt). Ein neuer Schlag
setzt zurück:
\begin{lstlisting}[style=py,caption={\code{flight\_ekf} -- Einfrieren und Reset}]
def on_point(self, msg):
    z = self._to_world(msg)                  # TF: optical -> world
    if self.frozen:                          # Schaetzung steht fest
        if z[1] > self.reset_y: self._reset()   # nur neuer Schlag setzt zurueck
        else: return                         # sonst: kein EKF-Rechnen mehr
    ...                                      # init / ekf_step / publish
    if self.x[1] <= self.freeze_y:           # Tischmitte erreicht
        self.frozen = True                   # -> Abfangpunkt unveraenderbar
\end{lstlisting}
\noindent Verifiziert: das Einfrieren erfolgt bei $y\approx0$; danach bleiben
Schätzung und Abfangpunkt identisch (keine weiteren Updates), ein Ball wieder bei
$y>1{,}0$ löst den Reset aus.

\begin{infobox}[title={Zwei EKF-Stolpersteine, die wir live trafen}]
\textbf{(1) Geschwindigkeit initialisieren:} mit Startgeschwindigkeit $0$ braucht
der EKF zu lange -- wir setzen $\mathbf v$ aus der \emph{Differenz der ersten
zwei Messungen} ($\mathbf v=(\mathbf z_1-\mathbf z_0)/\Delta t$). \textbf{(2) $Q$
strukturieren:} ein großes $Q$ auf \emph{allen} Zuständen macht die Schätzung
verrauscht; klein auf der Position (das Modell $\mathbf r\!+\!\mathbf v\,\Delta t$
ist exakt), moderat auf der Geschwindigkeit. Damit liegt die \emph{frühe}
Abfangpunkt-Vorhersage (nach $\sim$16 Messungen) bei wenigen bis $\sim$20\,cm und
verfeinert sich mit jeder weiteren Messung (frühe Reaktion, vgl.\ 26.1.8).
\end{infobox}

\begin{infobox}[title={Checkliste 23.5/24.6 -- in Schritt 4 beachtet}]
Mess-Frame korrekt behandelt (TF optical$\rightarrow$\code{world});
\code{use\_sim\_time} + \code{/clock}; Pose-/Punktfelder mit \code{float(\ldots)}
gecastet; \emph{kein} Custom-Interface (\code{PointStamped}, \code{nav\_msgs/Path},
\code{Marker}); Methoden neutral benannt (\code{ekf\_step}, nicht \code{handle}).
Zur Konsistenz nutzt auch \code{opponent\_serve} jetzt Luftwiderstand (gleiche
$k$), und die Schlagparameter sind so gewählt, dass der Ball \emph{über} das Netz
kommt.
\end{infobox}

\paragraph{Terminals und Befehle.}
\begin{lstlisting}[style=bash,caption={Terminal 1 -- alles starten (jetzt mit flight\_ekf)}]
cd ~/ros2_ws
colcon build --packages-select p2_tabletennis
source install/setup.bash
ros2 launch p2_tabletennis tabletennis.launch.py
\end{lstlisting}
\begin{lstlisting}[style=bash,caption={Terminal 2 -- Schaetzung und Abfangpunkt mitlesen}]
source ~/ros2_ws/install/setup.bash
ros2 topic echo /ball/estimate     # EKF-Schaetzung (im world-Frame)
ros2 topic echo /ball/intercept    # vorhergesagter Abfangpunkt
\end{lstlisting}
\noindent In RViz zeigt der gelbe Zylinder den vorhergesagten Abfangpunkt, die
blaue Linie die prädizierte (gebouncte) Bahn. Geht der Schlag ins Netz oder ins
Aus, verschwindet der Abfang-Marker (ungültig).
\begin{lstlisting}[style=bash,caption={\code{ros2\_ws}-Delta nach Schritt 4 (\# neu / \# angepasst)}]
src/p2_tabletennis/
|- p2_tabletennis/flight_ekf.py      # neu (EKF + Luftwiderstand + Auftreffpunkt)
|- p2_tabletennis/opponent_serve.py  # angepasst: Luftwiderstand + Schlag ueber Netz
|- launch/tabletennis.launch.py      # angepasst: flight_ekf statt 0.B-ball_tracker
|- config/tabletennis.rviz           # angepasst: Marker /ball/intercept_marker
`- setup.py / package.xml            # angepasst: flight_ekf, tf2_*-Abhaengigkeiten
\end{lstlisting}

\subsection{Schritt 5 -- Aufschlagregeln, Fehlschuss und Begrenzungswand}
Schritt~5 setzt die Tischtennis-Regeln aus 26.2.1.5 um und macht das Ergebnis
\emph{sichtbar}. Ein neuer Knoten \code{match\_referee} (Schiedsrichter) erkennt
die Aufpraller des Balls, prüft sie gegen die Aufschlag- bzw.\ Ballwechselregeln
und steuert die Visualisierung. Zusätzlich erhält \code{opponent\_serve} einen
echten \textbf{Aufschlag}-Modus -- denn die bisherigen Würfe waren Ballwechsel,
keine Aufschläge.

\paragraph{Regelkonformer Aufschlag im Gegnerknoten.} Der Roboter schlägt nie
selbst auf; der Aufschlag kommt immer vom Spieler. \code{opponent\_serve} kennt
daher zwei Modi (Services \code{/opponent/serve} und \code{/opponent/rally}) und
meldet jeden Schlag auf \code{/opponent/event}, womit der Schiedsrichter seine
Zustandsmaschine zurücksetzt. Die Aufschlagparameter sind so gewählt (offline
geprüft), dass der Ball zuerst auf der eigenen Seite ($y\approx+0{,}5$), dann über
das Netz ($z\approx0{,}97>0{,}9125$) auf die Roboterseite ($y\approx-0{,}5$) springt.
\begin{lstlisting}[style=py,caption={\code{opponent\_serve} -- Modus wählen und Schiedsrichter informieren}]
def _launch(self, mode):
    if mode == 'serve':   # 1. Aufprall eigene Seite, dann ueber Netz
        s, speed, az, el = self.serve_start, 3.0, -1.5708, 0.0873
    else:                 # Ballwechsel: direkt ueber Netz auf Roboterseite
        s, speed, az, el = self.start, 4.5, -1.5708, 0.10
    self.pos = [float(s[0]), float(s[1]), float(s[2])]
    self.vel = [speed*cos(el)*cos(az), speed*cos(el)*sin(az), speed*sin(el)]
    self.event_pub.publish(String(data=mode))   # Schiedsrichter zuruecksetzen
\end{lstlisting}

\paragraph{Aufprall-Erkennung aus den Messungen.} Der Schiedsrichter darf
\emph{nicht} die EKF-Schätzung verwenden -- die ist ab der Tischmitte eingefroren
(Schritt~4) und sähe den Aufprall auf der Roboterseite gar nicht. Stattdessen
transformiert er die Messungen \code{/ball/point} per TF nach \code{world} und
erkennt einen Aufprall als \emph{lokales $z$-Minimum} nahe der Tischhöhe (Vorzeichen
von $v_z$ kippt von $-$ nach $+$), innerhalb der Platte:
\begin{lstlisting}[style=py,caption={\code{match\_referee} -- Aufprall = $z$-Minimum auf Tischhöhe}]
def _check_bounce(self):
    (p0,t0),(p1,t1),(p2,t2) = self.buf          # letzte drei Messungen
    vz_old = (p1[2]-p0[2])/(t1-t0); vz_new = (p2[2]-p1[2])/(t2-t1)
    if not (vz_old < -0.2 and vz_new > 0.2):     return  # kein Umkehrpunkt
    if p1[2] > self.tz + 0.10:                   return  # nicht auf Tischhoehe
    if abs(p1[0]) > self.hw or abs(p1[1]) > self.hl:  return  # ausserhalb Platte
    if t1 - self.last_bounce_t < 0.15:           return  # Entprellung
    self.last_bounce_t = t1
    self._on_bounce('player' if p1[1] > 0 else 'robot')
\end{lstlisting}

\paragraph{Regel-Zustandsmaschine.} Aus der Reihenfolge der Aufprall-Seiten ergibt
sich der Verdikt. Ein Netztreffer (Bahn schneidet $y=0$ unter $0{,}9125\,$m) wird
zusätzlich beim Überqueren der Netzebene erkannt.
\begin{lstlisting}[style=py,caption={\code{match\_referee} -- Aufschlag-/Ballwechselregeln}]
def _on_bounce(self, side):
    self.bounces.append(side); n = len(self.bounces)
    if self.phase == 'SERVE':
        if n == 1 and side == 'robot':  self._fault('nicht auf Aufschlagseite')
        elif n == 2 and side == 'robot': self._valid_serve()  # -> Phase RALLY
        elif n == 2:                    self._fault('2x eigene Seite')
    else:  # RALLY
        if n == 1 and side == 'robot':  self._status('gueltiger Aufprall')
        elif n == 1:                    self._fault('falsche Seite')
        else:                           self._fault('Doppelaufprall')
\end{lstlisting}

\paragraph{Begrenzungswand und Auftreffpunkt in Gazebo.} Die Wand ist als
\emph{durchsichtiges, kollisionsfreies} Box-Visual am Plattenrand ($y=-1{,}30$,
der Abfangebene) in die Welt eingebaut -- der Ball fliegt hindurch. Da Gazebo zur
Laufzeit kein bequemes Umfärben über ROS erlaubt, nutzt der Schiedsrichter einen
Trick \emph{ohne} Materialeingriff: Es gibt \emph{zwei} Wände -- eine grüne
(\code{wall\_ok}) und eine rote (\code{wall\_fault}, sonst unter dem Boden
geparkt). Bei einem Fehlschuss werden sie über dieselbe \code{set\_pose}-Brücke
wie der Ball (Schritt~3) getauscht; ein neuer Ballwechsel tauscht zurück. Der vom
EKF geschätzte Auftreffpunkt (\code{/ball/intercept}) wird als flache Scheibe live
auf die Wand gesetzt.
\begin{lstlisting}[style=py,caption={\code{match\_referee} -- Wand umfärben \& Auftreffpunkt setzen (set\_pose)}]
def _set_wall(self, fault):                      # Tausch statt Umfaerben
    self._move('wall_ok',    0, self.wall_y, self.park_z if fault else self.wall_z)
    self._move('wall_fault', 0, self.wall_y, self.wall_z if fault else self.park_z)

def on_intercept(self, msg):                     # EKF-Schaetzung auf die Wand
    if self.fault: return
    self._move('impact_marker', msg.point.x, self.wall_y, msg.point.z,
               qx=0.7071, qw=0.7071)             # roll 90: Scheibe liegt flach
\end{lstlisting}

\begin{warnbox}[title={Schiedsrichter liest Messungen, nicht die EKF-Schätzung}]
Die Aufprall-Erkennung läuft auf \code{/ball/point} (per TF nach \code{world}),
\emph{nicht} auf \code{/ball/estimate}. Grund: Der EKF friert seine Schätzung an
der Tischmitte ein (Schritt~4) und liefert für die Roboterseite -- wo der
entscheidende zweite Aufprall liegt -- gar keine Updates mehr. Die rohe Messung
sieht die ganze Bahn.
\end{warnbox}

\begin{infobox}[title={Checkliste 23.5/24.6 -- in Schritt 5 beachtet}]
Mess-Frame korrekt (TF optical$\rightarrow$\code{world}); \code{use\_sim\_time} +
\code{/clock} an allen Knoten; \code{set\_pose} wird erst nach
\code{service\_is\_ready()} aufgerufen (kein Blockieren); Pose-Felder mit
\code{float(\ldots)} gecastet; \emph{keine} Custom-Interfaces (\code{PointStamped},
\code{std\_msgs/String}, \code{ros\_gz\_interfaces/SetEntityPose}); Knoten und
Methoden neutral benannt (\code{match\_referee}, \code{process\_world\_point}).
Die Wand ist \emph{ohne Kollision} modelliert -- sonst würde der Ball an ihr
abprallen statt hindurchzufliegen.
\end{infobox}

\paragraph{Terminals und Befehle.}
\begin{lstlisting}[style=bash,caption={Terminal 1 -- alles starten (jetzt mit Schiedsrichter + Wand)}]
cd ~/ros2_ws
colcon build --packages-select p2_tabletennis
source install/setup.bash
ros2 launch p2_tabletennis tabletennis.launch.py
\end{lstlisting}
\begin{lstlisting}[style=bash,caption={Terminal 2 -- Aufschlag/Ballwechsel auslösen, Verdikt mitlesen}]
source ~/ros2_ws/install/setup.bash
ros2 service call /opponent/serve std_srvs/srv/Trigger   # regelkonformer Aufschlag
ros2 service call /opponent/rally std_srvs/srv/Trigger   # Ballwechsel-Schlag
ros2 topic echo /match/status                            # bereit / gueltig / FEHLSCHUSS
\end{lstlisting}
\noindent In Gazebo zeigt die Scheibe auf der durchsichtigen Wand den
EKF-Auftreffpunkt; der reale Ball durchbricht die Wand daneben -- die Differenz ist
die Modellgüte. Bei einem ungültigen Aufschlag (z.\,B.\ \code{/opponent/serve} mit
einem Ballwechsel-artigen Wurf, oder ein Ball ins Netz) färbt sich die Wand rot,
beim nächsten Schlag wieder grün.
\begin{lstlisting}[style=bash,caption={\code{ros2\_ws}-Delta nach Schritt 5 (\# neu / \# angepasst)}]
src/p2_tabletennis/
|- p2_tabletennis/match_referee.py   # neu (Regeln + Fehlschuss + Wand-Visualisierung)
|- p2_tabletennis/opponent_serve.py  # angepasst: Aufschlag-/Ballwechsel-Modus + /opponent/event
|- worlds/tabletennis.sdf            # angepasst: wall_ok/wall_fault + impact_marker (visual-only)
|- launch/tabletennis.launch.py      # angepasst: match_referee gestartet
`- setup.py / package.xml            # angepasst: match_referee, std_msgs
\end{lstlisting}

\noindent Offline verifiziert (Regel-Logik mit synthetischen Bahnen, gleiches
Drag-Modell): gültiger Aufschlag (player$\rightarrow$robot) $\rightarrow$ Wechsel
auf Ballwechsel; erster Aufprall Roboterseite $\rightarrow$ Fehlschuss; zweiter
Aufprall erneut eigene Seite $\rightarrow$ Fehlschuss; Aufschlag ins Netz
$\rightarrow$ Fehlschuss (Netztreffer); gültiger Ballwechsel (ein Aufprall
Roboterseite) $\rightarrow$ kein Fehlschuss.

\subsection{Schritt 6 -- Roboteranbindung}
\emph{[Dokumentation folgt.]} Frühe Abfangbewegung (iterative Verfeinerung),
inverse Kinematik des Portals und Motorsteuerung über den
\code{ros2\_control}-Stack.
